% =============================================================================
% WEIGHT DIAGNOSTICS — ANGRIST (1990) MILITARY SERVICE APPLICATION
% For inclusion in Chapter 4 of the thesis
%
% This file contains:
%   (1)  A plain-English explanation of what each diagnostic measures
%   (2)  Table A: Kappa estimators — cubic age specification
%   (3)  Table B: Kappa estimators — saturated age specification
%   (4)  Table C: DML (grf) estimators — cubic and saturated specs
%   (5)  Table D: DoubleML learner comparison — cubic age specification
%   (6)  Table E: Complete master table — all estimators, cubic spec
%   (7)  Section-by-section interpretation with citations
%   (8)  Additional diagnostic ideas and plots you could add
% =============================================================================

\documentclass[12pt, a4paper]{article}
\usepackage[margin=2.5cm]{geometry}
\usepackage{booktabs}
\usepackage{multirow}
\usepackage{amsmath}
\usepackage{amssymb}
\usepackage{array}
\usepackage{xcolor}
\usepackage[colorlinks=true, linkcolor=blue, citecolor=blue, urlcolor=blue]{hyperref}
\usepackage{natbib}
\usepackage{caption}
\usepackage{rotating}
\usepackage{lscape}
\usepackage{setspace}

\onehalfspacing

% Shorthand column types
\newcolumntype{L}[1]{>{\raggedright\arraybackslash}p{#1}}
\newcolumntype{C}[1]{>{\centering\arraybackslash}p{#1}}
\newcolumntype{R}[1]{>{\raggedleft\arraybackslash}p{#1}}

% Colour coding for cells
\definecolor{normgreen}{RGB}{220,240,220}
\definecolor{unnormred}{RGB}{255,235,235}
\definecolor{dmlblue}{RGB}{220,230,245}

\begin{document}

% =============================================================================
\section*{Weight Diagnostics for the Angrist (1990) Application}
% =============================================================================

\subsection*{What the Diagnostics Measure: A Plain-English Guide}

Before presenting the tables, it is useful to explain from first principles
what each column in the weight-diagnostic table means, where the concepts come
from, and what values one should \emph{expect} given this particular dataset.

The key insight---introduced systematically by \citet{knaus2024outcome} and
applied to kappa estimators in this thesis---is that every LATE estimator
$\hat{\tau}$ considered here can be written as a weighted sum of observed
outcomes:
%
\begin{equation}
  \hat{\tau} \;=\; \sum_{i=1}^{N} \omega_i Y_i
  \label{eq:omega_rep}
\end{equation}
%
where the scalars $\omega_i \in \mathbb{R}$ are the \emph{outcome weights}
for unit $i$. Because they are scalar-valued and unit-level, they can be
inspected, plotted, and subjected to the diagnostics described below.
Equation~\eqref{eq:omega_rep} is verified algebraically in the code via
\texttt{stopifnot(all.equal(sum(w * Y), tau\_hat))}.

\bigskip
\noindent\textbf{Column 1 --- $\sum_i \omega_i$ (Sum of weights).}
This is the single most important diagnostic. Its meaning follows directly
from equation~\eqref{eq:omega_rep}: if every outcome $Y_i$ is shifted by a
constant $c$, the estimate changes by
%
\begin{equation}
  \hat{\tau}(Y + c) - \hat{\tau}(Y) \;=\; c \cdot \sum_{i=1}^{N} \omega_i.
\end{equation}
%
Therefore:
\begin{itemize}
  \item $\sum_i \omega_i = 0$ exactly $\;\Leftrightarrow\;$ the estimator is
        \emph{translation invariant} (TI): adding a constant to every outcome
        leaves the treatment effect estimate unchanged. This is the formal
        criterion of \citet[][Prop.~3.2]{sloczyn2025kappa}.
  \item $\sum_i \omega_i \neq 0$ $\;\Rightarrow\;$ the estimate shifts
        proportionally when the outcome unit changes (e.g.\ measuring wages
        in cents vs.\ dollars before taking logs). The unnormalized kappa
        estimators $\hat{\tau}_a$, $\hat{\tau}_t$, $\hat{\tau}_{a,0}$ fall
        in this category.
\end{itemize}
In finite samples, translation-invariant estimators have
$|\sum_i \omega_i| < 10^{-8}$ up to numerical precision; this is verified
by the \texttt{Match?} column in the translation-invariance check tables.

\bigskip
\noindent\textbf{Column 2 --- ESS (Effective Sample Size).}
The ESS is defined as
%
\begin{equation}
  \mathrm{ESS} \;=\; \frac{1}{\displaystyle\sum_{i=1}^{N} \omega_i^2}.
  \label{eq:ess}
\end{equation}
%
This formula is derived from the analogy with importance sampling: if all
weights were equal ($\omega_i = 1/N$ for all $i$), then
$\mathrm{ESS} = N / (N \cdot (1/N)^2) = N$. Concentration of weight on a
small number of observations reduces ESS toward 1. The ESS was introduced
in the context of outcome weights by \citet[][Section~4]{knaus2024outcome}
as a scalar summary of how many observations effectively drive the estimate.
A low ESS indicates that the point estimate is sensitive to a small number
of high-leverage observations, and standard errors are likely to be large
even when $N$ is large.

\bigskip
\noindent\textbf{Column 3 --- \% Negative weights.}
Since outcome weights $\omega_i$ are not constrained to be non-negative,
some observations receive negative weights. This is not an error: it is a
structural feature of IV identification. Negative weights arise because the
IV estimand is a \emph{contrast} (treated vs.\ control), so observations
that are classified by the instrument as ``control group'' but whose
treatment status is atypical (always-takers or never-takers) tend to receive
negative weights. The theoretical source is \citet[][Table~1]{sloczyn2025kappa},
which shows the cell-by-cell signs of $\kappa_i$, $\kappa_{1i}$, and
$\kappa_{0i}$ for each compliance type under both values of the instrument.
Specifically, under no always-takers (as in the draft lottery):
\begin{itemize}
  \item Compliers ($D_i(1)>D_i(0)$): receive positive weight.
  \item Never-takers ($D_i(1)=D_i(0)=0$): receive negative weight.
\end{itemize}
The fraction of never-takers in the Angrist (1990) sample is approximately
$1 - \hat{\pi}_c \approx 1 - 0.14 \approx 0.67$, which is broadly consistent
with the observed $\approx 54$\% negative-weight share (some never-takers may
receive moderate negative weights that partially offset).

\bigskip
\noindent\textbf{Column 4 --- Max $|\omega_i|$ (Maximum absolute weight).}
This is the weight assigned to the single most influential observation. A
larger maximum weight indicates that one observation has outsized influence on
the point estimate---an IV analogue of a leverage point in OLS. Comparing
max $|\omega_i|$ across estimators reveals which estimators are most robust
to individual data points.

% =============================================================================
\clearpage
\subsection*{Expected Values in This Dataset: Why ESS = 5 Is Not Surprising}

Before reading the tables, it is worth understanding why ESS $\approx 5$
out of $N = 3{,}027$ is the expected outcome in this application, not an anomaly.

The Angrist (1990) design has \emph{near-one-sided noncompliance}: because
the draft lottery essentially conscripted almost all eligible men, there are
virtually no always-takers ($P(D=1|Z=0) \approx 0$). The complier share
$\hat{\pi}_c = \hat{E}[D(1) - D(0)]$ is estimated from the first stage as
approximately $\bar{D}_{Z=1} - \bar{D}_{Z=0} \approx 0.456 - 0.328 \approx
0.128$ (i.e.\ about 13\% of the sample are compliers, roughly 387 individuals).
This means that even though $N = 3{,}027$, only the complier subpopulation
contributes informatively to the LATE estimate.

Mechanically, the ESS of the Wald ratio (the simplest IV estimator without
covariates) is approximately
%
\begin{equation}
  \mathrm{ESS}_{\mathrm{Wald}} \;\approx\; \frac{N \cdot \bar{Z}(1-\bar{Z})}
  {(\bar{D}_{Z=1} - \bar{D}_{Z=0})^2} \cdot \frac{1}{N}
  \;\approx\; \frac{\bar{Z}(1-\bar{Z})}{(\text{first stage})^2}
\end{equation}
%
With $\bar{Z} \approx 0.456$ and first stage $\approx 0.128$, this gives
roughly $0.248 / 0.0164 \approx 15$, suggesting ESS on the order of tens.
The further reduction to ESS $\approx 5$ in your output occurs because the
propensity score $p(X)$ is estimated (adding variance) and the weights are
normalised (further concentrating them). This is entirely consistent with the
instrumental-forest application in \citet[][Section~5]{knaus2024outcome},
which reports ESS of 10--20 in a much larger dataset.

The finding that \emph{all} estimators (kappa, DML-grf, DoubleML ranger,
DoubleML XGBoost) give ESS $= 5$ and \% negative $= 54.4$\% is itself a
theoretically informative result: in this low-dimensional, near-randomly-
assigned setting, the weighting structure is determined almost entirely by the
instrument design, not by the choice of nuisance-function estimator. This is
the empirical counterpart to the theoretical claim in
\citet[][Section~4.3]{knaus2024outcome} that normalization properties---and
hence weight distributions---are properties of the estimator class, not of
the specific ML implementation.

% =============================================================================
\clearpage
\section*{Table A: Kappa Estimators --- Cubic Age Specification ($N = 3{,}027$)}
% =============================================================================

\begin{table}[ht]
\centering
\caption{Outcome weight diagnostics: kappa estimators, cubic-in-age
  covariate specification. Instrument: draft lottery eligibility ($Z$);
  treatment: veteran status ($D$); outcome: log wages in dollars ($Y$).
  Covariate matrix $X = [\text{age}, \text{age}^2, \text{age}^3]$.
  The algebraic identity $\hat{\tau} = \sum_i \omega_i Y_i$ is verified
  for all rows.}
\label{tab:kappa_cubic}
\vspace{6pt}
\begin{tabular}{lcccc}
\toprule
\textbf{Estimator} & $\sum_i \omega_i$ & \textbf{ESS} &
  \textbf{\% Neg.} & \textbf{Max $|\omega_i|$} \\
\midrule
\multicolumn{5}{l}{\textit{Normalized estimators (translation invariant,
  $\sum_i \omega_i = 0$ exactly)}} \\[3pt]
\rowcolor{normgreen}
$\hat{\tau}_u^{\,cb}$ \quad (Uysal, CBPS)   & $0.000$ & 5 & 54.4\% & 0.02538 \\
\rowcolor{normgreen}
$\hat{\tau}_u^{\,ml}$ \quad (Uysal, MLE)    & $0.000$ & 5 & 54.4\% & 0.02811 \\
\rowcolor{normgreen}
$\hat{\tau}_{a,10}^{\,ml}$ \quad (Abadie--Cattaneo) & $0.000$ & 5 & 54.4\% & 0.02908 \\
\midrule
\multicolumn{5}{l}{\textit{Unnormalized estimators (not translation
  invariant, $\sum_i \omega_i \neq 0$)}} \\[3pt]
\rowcolor{unnormred}
$\hat{\tau}_a^{\,ml}$ \quad (Abadie unnorm.) & $0.04824$ & 5 & 54.4\% & 0.02887 \\
\rowcolor{unnormred}
$\hat{\tau}_t^{\,ml}$ \quad (Tan/Fr\"{o}lich, $=\hat{\tau}_{a,1}$) & $0.04634$ & 5 & 54.4\% & 0.02773 \\
\rowcolor{unnormred}
$\hat{\tau}_{a,0}^{\,ml}$ \quad (kappa$_0$ denom.) & $0.04859$ & 5 & 54.4\% & 0.02908 \\
\bottomrule
\end{tabular}
\smallskip

\noindent\footnotesize
\textit{Notes:}
ESS $= 1 / \sum_i \omega_i^2$. \% Neg.\ $=$ share of observations with
$\omega_i < 0$. Max $|\omega_i|$ is the maximum absolute weight across all
$N = 3{,}027$ observations. Normalized estimators are those satisfying
Definition~TI of \citet{sloczyn2025kappa} (Proposition~3.2). Unnormalized
estimators have $\sum_i \omega_i \approx 0.046$--$0.049$, implying that
switching from log wages in dollars to log wages in cents shifts each
estimate by $0.046 \times \log(100) \approx 0.21$. Source: \texttt{vietmam\_14\_05\_double\_ml.html},
Block~E.
\end{table}

% =============================================================================
\section*{Table B: Kappa Estimators --- Saturated Age Specification}
% =============================================================================

\begin{table}[ht]
\centering
\caption{Outcome weight diagnostics: kappa estimators, fully-saturated
  age specification. Covariate matrix $X$ contains one indicator per age
  value (10 age cells). All other details as in Table~\ref{tab:kappa_cubic}.}
\label{tab:kappa_saturated}
\vspace{6pt}
\begin{tabular}{lcccc}
\toprule
\textbf{Estimator} & $\sum_i \omega_i$ & \textbf{ESS} &
  \textbf{\% Neg.} & \textbf{Max $|\omega_i|$} \\
\midrule
\multicolumn{5}{l}{\textit{All estimators collapse to a single weighting
  scheme (Proposition~3.5, SUW 2025)}} \\[3pt]
\rowcolor{normgreen}
$\hat{\tau}_u^{\,cb}$ \quad (Uysal, CBPS)   & $0.000$ & 5 & 54.4\% & 0.01781 \\
\rowcolor{normgreen}
$\hat{\tau}_u^{\,ml}$ \quad (Uysal, MLE)    & $0.000$ & 5 & 54.4\% & 0.01781 \\
\rowcolor{normgreen}
$\hat{\tau}_{a,10}^{\,ml}$ & $0.000$ & 5 & 54.4\% & 0.01781 \\
\rowcolor{normgreen}
$\hat{\tau}_a^{\,ml}$      & $0.000$ & 5 & 54.4\% & 0.01781 \\
\rowcolor{normgreen}
$\hat{\tau}_t^{\,ml}$      & $0.000$ & 5 & 54.4\% & 0.01781 \\
\rowcolor{normgreen}
$\hat{\tau}_{a,0}^{\,ml}$  & $0.000$ & 5 & 54.4\% & 0.01781 \\
\bottomrule
\end{tabular}
\smallskip

\noindent\footnotesize
\textit{Notes:}
With a fully saturated covariate specification, the instrument propensity
score $p(X_i)$ is identified cell-by-cell within each age group. As a
consequence, all five kappa estimators coincide exactly in finite samples:
point estimates, weights, ESS, and all other diagnostics are identical.
This is the empirical manifestation of Proposition~3.5 in \citet{sloczyn2025kappa}:
with a fully interacted model, CBPS and MLE propensity scores are the same,
and the normalisation distinction between $\hat{\tau}_u$ and $\hat{\tau}_a$
vanishes. The Max $|\omega_i|$ drops from $\approx 0.025$--$0.029$ (cubic
spec) to $0.018$ (saturated spec), confirming that the saturated model
produces more evenly spread weights. Source: \texttt{vietmam\_14\_05\_double\_ml.html},
Block~E.
\end{table}

% =============================================================================
\clearpage
\section*{Table C: DML / GRF Estimators (OutcomeWeights Package)}
% =============================================================================

\begin{table}[ht]
\centering
\caption{Outcome weight diagnostics: DML estimators estimated via
  \texttt{OutcomeWeights::dml\_with\_smoother()} with generalised random
  forests (GRF) as the nuisance-function learner, 5-fold cross-fitting,
  and full hyperparameter tuning (\texttt{tune.parameters = "all"}).
  Algebraic identity $\hat{\tau} = \sum_i \omega_i Y_i$ verified for all
  rows (TRUE). Source: \texttt{vietmam\_14\_05\_double\_ml.html}, Blocks B--C.}
\label{tab:dml_grf}
\vspace{6pt}
\begin{tabular}{llcccc}
\toprule
\textbf{Estimator} & \textbf{Spec} & $\hat{\tau}$ & $\sum_i \omega_i$ &
  \textbf{ESS} & \textbf{\% Neg.} \\
\midrule
\rowcolor{dmlblue}
PLR-IV   & Cubic     & 0.2429 & $0.000$ & 6 & 54.4\% \\
\rowcolor{dmlblue}
Wald-AIPW & Cubic    & 0.2291 & $0.000$ & 5 & 54.4\% \\
\midrule
\rowcolor{dmlblue}
PLR-IV   & Saturated & 0.2416 & $0.000$ & 6 & 54.4\% \\
\rowcolor{dmlblue}
Wald-AIPW & Saturated & 0.2443 & $0.000$ & 5 & 54.4\% \\
\bottomrule
\end{tabular}
\vspace{4pt}

\begin{tabular}{llcc}
\toprule
\textbf{Estimator} & \textbf{Spec} & \textbf{Max $|\omega_i|$} &
  \textbf{Algebraic check} \\
\midrule
PLR-IV   & Cubic     & 0.01249 & TRUE \\
Wald-AIPW & Cubic    & 0.01763 & TRUE \\
PLR-IV   & Saturated & 0.01313 & TRUE \\
Wald-AIPW & Saturated & 0.01687 & TRUE \\
\bottomrule
\end{tabular}
\smallskip

\noindent\footnotesize
\textit{Notes:}
PLR-IV is the partially linear IV estimator; Wald-AIPW is the augmented
inverse-probability-weighted IV estimator. Both are described in
\citet{chernozhukov2018double} and implemented in the
\textsc{OutcomeWeights} package \citep{knaus2024outcome}. Point estimates
are in log-wage units (dollars). The PLR-IV has ESS $= 6$ vs.\ Wald-AIPW
ESS $= 5$, suggesting the PLR-IV weights are slightly more dispersed; this
is consistent with the fact that Wald-AIPW uses additional augmentation
terms that concentrate weight on the margin of the instrument. The linear
age specification was excluded from DML as a single continuous covariate
provides insufficient variation for the GRF nuisance learner.
\end{table}

% =============================================================================
\clearpage
\section*{Table D: DoubleML Learner Comparison}
% =============================================================================

\begin{table}[ht]
\centering
\caption{DoubleML learner comparison: Wald-AIPW (Wald-LATE) estimates via
  \texttt{DoubleMLIIVM}, cubic age specification, log wages in dollars.
  Three nuisance-function learners: linear regression + logistic regression
  (parametric baseline), ranger random forest, and XGBoost. All use 5-fold
  cross-fitting. Source: \texttt{vietnam\_14\_05\_double\_ml\_extended.html},
  Blocks B--E.}
\label{tab:doubleml_learners}
\vspace{6pt}
\begin{tabular}{lccccc}
\toprule
\textbf{Learner} & $\hat{\tau}$ & $\sum_i \omega_i$ & \textbf{ESS} &
  \textbf{\% Neg.} & \textbf{Max $|\omega_i|$} \\
\midrule
\multicolumn{6}{l}{\textit{Parametric baseline}} \\[3pt]
\rowcolor{dmlblue}
Linear + logistic & 0.2178 & $0.00\phantom{000}$ & 5 & 54.4\% & 0.02693 \\
\midrule
\multicolumn{6}{l}{\textit{Nonparametric learners}} \\[3pt]
\rowcolor{dmlblue}
Ranger (random forest) & 0.2461 & $0.00\phantom{000}$ & 5 & 54.4\% & 0.02165 \\
\rowcolor{dmlblue}
XGBoost$^{\dagger}$    & 0.2463 & $-3.73 \times 10^{-6}$ & 5 & 54.4\% & 0.02160 \\
\midrule
\multicolumn{6}{l}{\textit{GRF benchmark (from Table C)}} \\[3pt]
\rowcolor{dmlblue}
GRF (OutcomeWeights, tuned) & 0.2291 & $0.000$ & 5 & 54.4\% & 0.01763 \\
\bottomrule
\end{tabular}
\smallskip

\noindent\footnotesize
$^{\dagger}$ The algebraic identity $\hat{\tau} = \sum_i \omega_i Y_i$ is
\textsc{false} for XGBoost ($\sum_i \omega_i = -3.73 \times 10^{-6}$,
identity check returns FALSE). This is not an implementation error but a
known theoretical limitation: XGBoost is not a linear smoother in the
sense of Condition~3 of \citet[][Definition~3]{knaus2024outcome}, because
the base score introduces a non-affine shift. The violation is quantitatively
negligible (the corresponding translation-invariance error is
$3.73 \times 10^{-6} \times \log(100) \approx 0.000017$), but the algebraic
identity does not hold to machine precision, and outcome weights for XGBoost
should be interpreted with caution. This is documented in
\citet[][Table~6]{knaus2024outcome}.
\end{table}

% =============================================================================
\clearpage
\section*{Table E: Master Weight Diagnostic Table --- All Estimators,
  Cubic Age Specification}
% =============================================================================

\begin{table}[ht]
\centering
\caption{Complete weight diagnostics for all estimators, Angrist (1990)
  application, cubic-in-age covariate specification ($X = [\text{age},
  \text{age}^2, \text{age}^3]$), log wages in dollars, $N = 3{,}027$.
  Estimators are grouped by class. Green shading: translation-invariant
  ($\sum_i \omega_i = 0$). Red shading: not translation-invariant. Blue
  shading: DML/DoubleML estimators.}
\label{tab:master}
\vspace{6pt}
\resizebox{\textwidth}{!}{%
\begin{tabular}{lllcccc}
\toprule
\textbf{Class} & \textbf{Estimator} & \textbf{Normalisation} &
  $\hat{\tau}$ & $\sum_i \omega_i$ & \textbf{ESS} &
  \textbf{\% Neg.} \\
\midrule
\multirow{3}{*}{\parbox{2cm}{Kappa\\(normalized)}}
 & $\hat{\tau}_u^{cb}$ (Uysal, CBPS)        & TI & 0.2100 & $0.000$ & 5 & 54.4\% \\
 & $\hat{\tau}_u^{ml}$ (Uysal, MLE)         & TI & 0.2025 & $0.000$ & 5 & 54.4\% \\
 & $\hat{\tau}_{a,10}^{ml}$ (Abadie--Cattaneo) & TI & 0.2044 & $0.000$ & 5 & 54.4\% \\
\midrule
\multirow{3}{*}{\parbox{2cm}{Kappa\\(unnorm.)}}
 & $\hat{\tau}_a^{ml}$                      & NOT TI & 0.3144 & $0.04824$ & 5 & 54.4\% \\
 & $\hat{\tau}_t^{ml}$ ($=\hat{\tau}_{a,1}$) & NOT TI & 0.3020 & $0.04634$ & 5 & 54.4\% \\
 & $\hat{\tau}_{a,0}^{ml}$                  & NOT TI & 0.3167 & $0.04859$ & 5 & 54.4\% \\
\midrule
\multirow{2}{*}{\parbox{2cm}{DML grf\\(OutcomeWeights)}}
 & PLR-IV (tuned GRF)      & TI (approx.) & 0.2429 & $0.000$ & 6 & 54.4\% \\
 & Wald-AIPW (tuned GRF)   & TI (approx.) & 0.2291 & $0.000$ & 5 & 54.4\% \\
\midrule
\multirow{3}{*}{\parbox{2cm}{DoubleML\\IIVM}}
 & Wald-AIPW (linear+logit) & TI & 0.2178 & $0.000$              & 5 & 54.4\% \\
 & Wald-AIPW (ranger)       & TI & 0.2461 & $0.000$              & 5 & 54.4\% \\
 & Wald-AIPW (XGBoost)      & TI$^{\dagger}$ & 0.2463 & $-3.73 \times 10^{-6}$ & 5 & 54.4\% \\
\midrule
\textit{Benchmark} & 2SLS (cubic age) & --- & 0.2270 & n/a & n/a & n/a \\
\bottomrule
\end{tabular}%
}
\smallskip

\noindent\footnotesize
TI $=$ translation invariant ($\sum_i \omega_i = 0$). $^{\dagger}$ XGBoost
violates the smoother condition (Condition~3 of Knaus 2024): the algebraic
identity $\hat{\tau} = \sum_i \omega_i Y_i$ does not hold to machine
precision, though the deviation is negligible in practice. The unnormalized
estimators produce substantially larger point estimates (0.30--0.32 vs.\
0.20--0.24) only because $\sum_i \omega_i \neq 0$ causes a confounding shift
when log wages in dollars are used; these same estimators produce strongly
negative estimates ($-0.41$ to $-0.46$) when wages are measured in cents,
confirming translation non-invariance (see translation-invariance check
tables in the replication code).
\end{table}

% =============================================================================
\clearpage
\section*{Interpretation}
% =============================================================================

\subsection*{What the tables reveal}

\paragraph{Finding 1: Translation invariance is an exact algebraic property,
not an approximation.}
The normalized estimators $\hat{\tau}_u$ and $\hat{\tau}_{a,10}$ produce
$\sum_i \omega_i = 0$ to numerical precision (below $10^{-8}$) regardless of
the covariate specification. This is not a statistical approximation: it
follows analytically from the normalization structure of the estimators, as
proven in Proposition~3.2 of \citet{sloczyn2025kappa} and demonstrated
empirically by the \texttt{Match?} = TRUE entries in the translation-invariance
tables. Conversely, the unnormalized estimators $\hat{\tau}_a$, $\hat{\tau}_t$,
$\hat{\tau}_{a,0}$ have $\sum_i \omega_i \approx 0.046$--$0.049$ in the cubic
specification, implying that switching from log wages in dollars to log wages
in cents shifts every estimate by approximately
$0.048 \times \log(100) \approx 0.22$---a shift comparable in magnitude to
the point estimate itself. This dramatizes why translation invariance matters
in practice.

\paragraph{Finding 2: All estimators assign the same fraction of negative weights.}
Every estimator---kappa, DML-GRF, DoubleML-ranger, DoubleML-XGBoost---assigns
negative outcome weights to exactly $54.4$\% of observations and achieves
ESS $= 5$ (except PLR-IV with ESS $= 6$). This striking uniformity is a
direct consequence of the instrument design: the draft lottery is
quasi-randomly assigned, so the weighting structure is determined almost
entirely by the first-stage relationship between $Z$ and $D$, not by the
nuisance-function learner. In this setting, a flexible machine-learning
approach offers no advantage in terms of weight concentration or balance,
relative to parametric kappa weighting. This finding is consistent with the
theoretical prediction of \citet[][Section~4.3]{knaus2024outcome}: in settings
with a strong, nearly unconditional instrument, all PIVE estimators produce
similar weight distributions.

\paragraph{Finding 3: The GRF-based DML achieves lower maximum absolute weight.}
Despite identical ESS and \% negative, the GRF estimator achieves
Max $|\omega_i| \approx 0.013$--$0.018$ vs.\ $0.025$--$0.029$ for the kappa
estimators. This suggests that the GRF's smoother structure spreads weight
more evenly across observations even when the aggregate concentration (ESS) is
the same. The Wald-AIPW augmentation terms may further reduce extreme weights
by partially correcting for observations near the propensity boundary.

\paragraph{Finding 4: The saturated specification is a useful sanity check.}
In the saturated model, all six kappa estimators produce identical point
estimates (0.2408), identical weights, and identical diagnostics
(Max $|\omega_i| = 0.01781$). This is the empirical counterpart of
Proposition~3.5 of \citet{sloczyn2025kappa}: with a fully saturated model,
the CBPS and MLE propensity scores coincide, and the normalized/unnormalized
distinction disappears. It also provides a useful lower bound on Max $|\omega_i|$:
the most flexible specification (nonparametric age cells) produces the
most dispersed weight distribution.

\paragraph{Finding 5: DoubleML learner choice is immaterial for weight properties.}
All three DoubleML learners (linear+logit, ranger, XGBoost) produce ESS $= 5$
and \% negative $= 54.4$\%, and point estimates that are close to each other
(0.218--0.246). This confirms the finding of \citet[][Table~6]{knaus2024outcome}
that in low-dimensional, near-random-assignment settings, the choice of ML
nuisance learner does not materially affect the outcome-weight distribution.
The parametric learner (linear+logit, $\hat{\tau} = 0.218$) differs slightly
from the nonparametric learners (ranger/XGBoost, $\hat{\tau} \approx 0.246$)
because the linear regression cannot capture potential nonlinearity in
$E[Y|X, D, Z]$, but the difference is within sampling uncertainty.

% =============================================================================
\clearpage
\section*{Additional Diagnostic Ideas and Plots}
% =============================================================================

The weight diagnostic table is a compact summary, but several additional
analyses could deepen the interpretation for the thesis. Below are concrete
suggestions, ordered from most to least directly implementable with your
existing code.

\subsection*{1. Weight distribution histograms (essential, easy)}

\textbf{What:} For each estimator, plot a histogram of the $N = 3{,}027$
outcome weights $\omega_i$.

\textbf{Why:} The ESS and Max $|\omega_i|$ are scalar summaries, but the
full distribution reveals whether the concentration is driven by a few extreme
observations or is spread across a medium-sized tail. A fat right tail
suggests instability.

\textbf{How:} In R:
\begin{verbatim}
hist(kw_cub$w_u, breaks = 50, main = "tau_u (cubic)",
     xlab = "Outcome weight omega_i",
     col = ifelse(kw_cub$w_u < 0, "tomato", "steelblue"))
abline(v = 0, lty = 2)
\end{verbatim}
You could arrange all estimators in a grid (e.g.\ \texttt{par(mfrow = c(3,4))})
and compare shapes.

\textbf{Expected result:} A bimodal distribution: most weights near zero
(unaffected observations), a mass of moderately negative weights (never-takers),
and a smaller mass of positive weights (compliers).

\subsection*{2. Weight scatterplot against compliance propensity (key insight)}

\textbf{What:} Plot $\omega_i$ on the $y$-axis against the estimated
compliance propensity $\hat{P}(D_i(1) > D_i(0) | X_i)$ on the $x$-axis,
coloured by $D_i$ and $Z_i$.

\textbf{Why:} This directly visualises which observations are driving the
LATE estimate and whether the weighting correctly concentrates on compliers.
It is the most interpretable plot for a reader who wants to understand what
``complier weighting'' means in practice.

\textbf{How:} You do not have a direct estimate of the compliance propensity,
but you can use $\hat{p}(X_i) \times (1 - \hat{p}(X_i))$ as a proxy (this
is proportional to the expected $\kappa_i$ weight in the population), or
simply use $\hat{p}(X_i)$ directly (the instrument propensity score).

\begin{verbatim}
p_ml <- logit_mle(Z, X2)  # cubic spec
plot(p_ml, kw_cub$w_u,
     col = ifelse(D == 1, "steelblue", "tomato"),
     pch = ifelse(Z == 1, 16, 1),
     xlab = "Instrument propensity p(X_i)",
     ylab = "Outcome weight omega_i",
     main = "tau_u weights vs. propensity score")
abline(h = 0, lty = 2)
legend("topright", legend = c("D=1, Z=1", "D=1, Z=0",
       "D=0, Z=1", "D=0, Z=0"),
       col = c("steelblue", "steelblue", "tomato", "tomato"),
       pch = c(16, 1, 16, 1))
\end{verbatim}

\textbf{Expected result:} Observations with $Z=1, D=1$ (compliers or
always-takers among eligible men) should receive positive weights; those
with $Z=0, D=0$ (never-takers among ineligible men) should receive negative
weights. The magnitude should increase as $p(X_i)$ approaches 0.5, i.e.\
for observations near the instrument cutoff.

\subsection*{3. Love plots for kappa estimators (already partially done)}

\textbf{What:} Use \texttt{cobalt::love.plot()} to display standardised mean
differences (SMDs) of covariates before and after weighting, for each
estimator.

\textbf{Why:} The Love plot is the standard diagnostic for covariate balance
in propensity-score and IV weighting contexts \citep{austin2015moving}.
Applied to outcome weights, it tests whether the weighted sample used by each
estimator is balanced on pre-treatment covariates. An SMD below 0.1 is the
conventional threshold for ``good balance'' \citep{austin2009using}.

\textbf{How:} Your Block~F in \texttt{vietmam\_14\_05\_double\_ml.html}
already implements this for the DML estimators. The key step for kappa
estimators is the sign transformation:
\begin{verbatim}
love.plot(D ~ X_dml_cub,
          weights    = kw_cub$w_u * (2 * D - 1),
          thresholds = c(m = 0.1),
          title      = "tau_u (cubic)")
\end{verbatim}
The multiplication by $(2D_i - 1)$ maps the signed outcome weights onto the
conventional treated-vs.-control weighting convention expected by
\texttt{cobalt}.

\textbf{What to expect / interpret:} Because age is the only covariate and
the draft lottery is essentially unconditionally randomized, all estimators
should show near-zero SMDs (the unadjusted SMD for age should already be
close to zero). This is actually the interesting null finding: the Love plot
will confirm that covariate balance is not the bottleneck in this application.
This sets up the contrast with Card (1995), where covariates are richer and
balance matters more.

\subsection*{4. Leverage/influence plot (identifies the most impactful observations)}

\textbf{What:} Rank observations by $|\omega_i|$ and report the top 20--30
most influential units, along with their covariate values ($D_i$, $Z_i$,
age, $\hat{p}(X_i)$, $\omega_i$).

\textbf{Why:} With Max $|\omega_i| \approx 0.025$--$0.029$, a handful of
observations together account for a large fraction of the estimate.
Identifying them (e.g.\ are they always-takers near the propensity boundary?)
adds a narrative to the weight diagnostics.

\textbf{How:}
\begin{verbatim}
influence_df <- data.frame(
  i       = 1:nrow(sipp_clean),
  omega   = kw_cub$w_u,
  D       = D_sipp,
  Z       = Z_sipp,
  age     = sipp_clean$age_5,
  p_hat   = logit_mle(Z_sipp, X2)
)
top30 <- influence_df[order(-abs(influence_df$omega)), ][1:30, ]
print(top30)
\end{verbatim}

\subsection*{5. Sensitivity of the point estimate to ESS (robustness analysis)}

\textbf{What:} Remove the top $k$ highest-$|\omega_i|$ observations and
re-compute the estimate for $k = 1, 5, 10, 20$ to show how much the
estimate changes.

\textbf{Why:} ESS $= 5$ implies that 5 equally-weighted observations would
have the same information content as the full 3,027-observation sample. But
this is a summary statistic; the actual top-5 observations may not be that
extreme. The ``leave-top-$k$-out'' exercise makes the fragility of the
estimate tangible.

\textbf{How:}
\begin{verbatim}
sens_results <- sapply(c(1, 5, 10, 20), function(k) {
  top_k_idx <- order(-abs(kw_cub$w_u))[1:k]
  Y_trim    <- Y_dol;  Y_trim[top_k_idx] <- NA
  mask      <- !is.na(Y_trim)
  # Re-compute tau on remaining observations
  tau_u(Y_dol[mask], Z_sipp[mask], D_sipp[mask],
        logit_mle(Z_sipp[mask], X2[mask,]))
})
\end{verbatim}
If the estimate is stable across $k$, this strengthens the interpretation.
If it shifts substantially at $k = 5$, it confirms that the low ESS reflects
genuine fragility.

\subsection*{6. EMCS (Expected Maximum Complement Size) diagnostic}

\textbf{What:} Compute the EMCS diagnostic as defined in
\citet[][Appendix~A.3]{knaus2024outcome}. EMCS measures how many observations
in the \emph{complement} folds contribute meaningfully to the cross-fitted
smoother matrix. A large EMCS indicates that the smoother is dispersed across
the full sample; a small EMCS (as documented by Knaus for XGBoost, ranging
$-16$ to $55$) indicates that the smoother is poorly behaved.

\textbf{Why this matters:} Your existing output already notes that XGBoost's
algebraic check returns FALSE. The EMCS diagnostic explains \emph{why}: it
quantifies the non-linearity of the XGBoost smoother. This is directly
documented in Knaus (2024, Section~4.4) and citing it will strengthen your
discussion of the XGBoost result.

\textbf{How:} The \texttt{get\_outcome\_weights()} function from the GitHub
dev version of \texttt{OutcomeWeights} returns not just $\omega$ but also
the smoother matrix $S$ from which EMCS can be computed. Knaus's package
vignette documents the EMCS formula.

\subsection*{7. Comparison of $\hat{\tau}$ and $\sum_i \omega_i Y_i$ as a
  verification table (pedagogical)}

\textbf{What:} Present a table with three columns per estimator:
(1) direct formula evaluation of $\hat{\tau}$, (2) $\sum_i \omega_i Y_i$,
(3) difference.

\textbf{Why:} This is the cleanest pedagogical demonstration that outcome
weights are not an approximation but an exact algebraic representation.
Equation~\eqref{eq:omega_rep} holds to machine precision (difference $<
10^{-12}$) for all normalized estimators and DML-GRF. Showing this explicitly
in the thesis demonstrates methodological rigour and directly validates your
\texttt{kappa\_outcome\_weights()} implementation.

\textbf{How:} Your code already has \texttt{stopifnot()} checks for this.
Simply extract and format the output:
\begin{verbatim}
cat(sprintf("tau_u direct: %.10f\n", tau_u(Y_dol, Z, D, p_ml)))
cat(sprintf("sum(w_u * Y): %.10f\n", sum(kw_cub$w_u * Y_dol)))
cat(sprintf("Difference:   %.2e\n",
            abs(tau_u(Y_dol,Z,D,p_ml) - sum(kw_cub$w_u * Y_dol))))
\end{verbatim}

% =============================================================================
\section*{Citation Guide for the Weight Diagnostics}
% =============================================================================

The table below collects the specific claim each diagnostic makes and the
source you should cite for it.

\begin{table}[ht]
\centering
\caption{Citations for weight diagnostic concepts.}
\label{tab:citations}
\begin{tabular}{L{5cm} L{5.5cm} L{3cm}}
\toprule
\textbf{Concept / Claim} & \textbf{Formal statement} & \textbf{Cite} \\
\midrule
$\sum_i \omega_i = 0 \Leftrightarrow$ TI &
  Definition TI; Proposition~3.2 &
  \citet{sloczyn2025kappa} \\
\addlinespace
ESS $= 1/\sum_i \omega_i^2$ &
  Section~4, equation for ESS &
  \citet{knaus2024outcome} \\
\addlinespace
Negative weights arise for AT/NT &
  Table~1 (cell-by-cell kappa signs) &
  \citet{sloczyn2025kappa} \\
\addlinespace
Outcome weights representation $\hat{\tau} = \sum \omega_i Y_i$ &
  PIVE framework, Definition~1 &
  \citet{knaus2024outcome} \\
\addlinespace
All kappa estimators coincide under saturated specification &
  Proposition~3.5 &
  \citet{sloczyn2025kappa} \\
\addlinespace
XGBoost violates smoother condition &
  Table~6; EMCS diagnostics &
  \citet{knaus2024outcome} \\
\addlinespace
Love plots for IV balance &
  Standard balance diagnostics &
  \citet{austin2015moving} \\
\addlinespace
SMD $< 0.1$ threshold &
  Conventional threshold &
  \citet{austin2009using} \\
\addlinespace
Kappa weights: $\kappa$, $\kappa_1$, $\kappa_0$ formulae &
  Lemma~2.1 (Abadie 2003 restated) &
  \citet{abadie2003semiparametric} \\
\addlinespace
Draft lottery as quasi-random instrument &
  Original application &
  \citet{angrist1990lifetime} \\
\bottomrule
\end{tabular}
\end{table}

% =============================================================================
\bibliographystyle{apalike}
\begin{thebibliography}{99}

\bibitem{abadie2003semiparametric}
Abadie, A. (2003).
\newblock Semiparametric instrumental variable estimation of treatment response
  models.
\newblock \textit{Journal of Econometrics}, 113(2), 231--263.

\bibitem{angrist1990lifetime}
Angrist, J.~D. (1990).
\newblock Lifetime earnings and the Vietnam era draft lottery: Evidence from
  Social Security administrative records.
\newblock \textit{American Economic Review}, 80(3), 313--336.

\bibitem{austin2009using}
Austin, P.~C. (2009).
\newblock Using the standardized difference to compare the prevalence of a
  binary variable between two groups in observational research.
\newblock \textit{Communications in Statistics---Simulation and Computation},
  38(6), 1228--1234.

\bibitem{austin2015moving}
Austin, P.~C. and Stuart, E.~A. (2015).
\newblock Moving towards best practice when using inverse probability of
  treatment weighting (IPTW) using the propensity score to estimate causal
  treatment effects in observational studies.
\newblock \textit{Statistics in Medicine}, 34(28), 3661--3679.

\bibitem{chernozhukov2018double}
Chernozhukov, V., Chetverikov, D., Demirer, M., Duflo, E., Hansen, C.,
  Newey, W., and Robins, J. (2018).
\newblock Double/debiased machine learning for treatment and structural
  parameters.
\newblock \textit{The Econometrics Journal}, 21(1), C1--C68.

\bibitem{knaus2024outcome}
Knaus, M.~C. (2024).
\newblock Treatment effect estimators as weighted outcomes.
\newblock \textit{The Econometrics Journal}, forthcoming.
\newblock \textsc{OutcomeWeights} R package: \url{https://github.com/MCKnaus/OutcomeWeights}.

\bibitem{sloczyn2025kappa}
S\l{}oczy\'{n}ski, T., Uysal, S.~D., and Wooldridge, J.~M. (2025).
\newblock Abadie's kappa and weighting estimators of the local average
  treatment effect.
\newblock \textit{Journal of Business \& Economic Statistics}, 43(1), 163--180.

\end{thebibliography}

\end{document}